%% Compiler: XeLaTeX (Overleaf: Menu -> Compiler -> XeLaTeX)
\documentclass[12pt]{article}
\usepackage[UTF8]{ctex}
\usepackage[margin=1in]{geometry}
\usepackage{booktabs,array,longtable,enumitem,hyperref}
\onehalfspacing
\newcommand{\path}[1]{\texttt{\detokenize{#1}}}
\title{Manuscript Outline}
\author{112700019 Paul}
\date{}
\begin{document}
\maketitle
This outline maps the R1--R4 draft folders to manuscript sections. Markers such as [R2] and [R3] identify source material that can be adapted directly.
\section*{Title}
\textit{Does the Relation, or Its Change Over Time, Predict Stock Returns? A Replication of Temporal Relational Ranking with a Dynamic-Relation Extension} (provisional; revise before finalizing).
\section*{Abstract}
One paragraph (150--250 words): question, method, and main findings. Write this last, after the R3/R4 results are final.
\section{Introduction}
Source: \path{reports/R1_topic_and_motivation/PLAN.md}. A draft exists and needs polishing into formal prose.
\begin{itemize}[leftmargin=*]
\item Motivation: why should stocks not be modeled as independent entities?
\item Feng et al. (2019), RSR: summarize the core method.
\item Research questions: RQ1 (replication) and RQ2 (dynamic-relations extension).
\item One-sentence summary of the contribution.
\end{itemize}
\section{Related Work}
Discuss HGTAN (hypergraph) and GCNET (graph convolution), and select two to four papers on supply-chain relations or cross-market information transmission.
\section{Data}
\textbf{[R2]} Source: \path{reports/R2_data_and_eda/eda.ipynb}.
\begin{itemize}[leftmargin=*]
\item Data source: 1,026 NASDAQ stocks in \path{data/2013-01-01/}, with 1,246 trading days.
\item Industry graph: 156 tickers (15\%) are ETFs/funds rather than individual stocks.
\item Describe the \path{-1234} missing-value code and sample size after ETF exclusion.
\end{itemize}
\section{Methods}
\textbf{[R3]} Source: \path{reports/R3_model_and_experiments/PROGRESS.md}.
\begin{itemize}[leftmargin=*]
\item Rank\_LSTM baseline without a relation graph.
\item RSR with a static industry graph and Temporal Graph Convolution.
\item Dynamic-relations extension using rolling-window correlations.
\end{itemize}
\section{Experiments and Results}
Report MSE, MRR-Top1, and simulated returns. Compare baseline, RSR, and dynamic-relations results using the same metrics. Current baseline results: Test MSE 0.000377, mrrt 0.049, and btl 1.02 (NASDAQ, 50 epochs). Other results require full runs.
\section{Discussion and Conclusion}
\textbf{[R4]} Compare the results with the original paper, explain differences in data period and market coverage, and state limitations. Address the syllabus prompt, ``Potential of your projects.'' If the extension proves useful, a next step could be a thesis or journal publication.
\section*{References}
Use Chicago style and include the papers listed in \path{references/}.
\section*{Writing Status}
Introduction, Related Work, Data, Methods, and baseline results have source material. RSR and dynamic-relations results require full training runs. Write the abstract last.
\end{document}